% =========================================================
\chapter{Lecture 6 : Applied cryptography II}
% =========================================================

\begin{rembox}[title=In a nutshell]
\textbf{What is this lecture about?} Lecture 5 only hid messages (confidentiality). Here we
fix the two holes it left: \textbf{authenticity} (am I really talking to Bob, and not to
Eve in the middle?) and \textbf{integrity} (was the message changed on the way?). With
public keys we get \textbf{public-key encryption} and \textbf{digital signatures}, made
fast thanks to \textbf{hash functions}; with a shared secret key we get \textbf{block
ciphers} and their \textbf{modes of operation} (ECB, CBC, CTR) and \textbf{MACs}. Finally
we combine encryption and MAC correctly (\textbf{Encrypt-then-MAC}). \\
\textbf{Why is it useful?} Every real secure channel (HTTPS, messaging apps, software
updates) needs all three properties at once. Knowing which tool gives which guarantee
--- and which naive combinations leak or can be tampered with --- is what lets you
choose the right standard primitive instead of building a broken one.
\end{rembox}

\section{Recap: where we stand}

\begin{rembox}[title=Recap of Lecture 5]
\begin{itemize}
  \item \textbf{Confidentiality}: information cannot be accessed by unauthorized
        parties. Eve observes the ciphertext $C = E_K(M)$, but is assumed \textbf{not to
        see inside} Alice's and Bob's devices (no access to $K$ or $M$).
  \item \textbf{Symmetric encryption}: encryption and decryption use the \textbf{same
        pre-shared secret key} $K$.
  \item \textbf{Diffie-Hellman}: with public $p, g$, Bob sends $B = g^b \bmod p$, Alice
        sends $A = g^a \bmod p$; both compute $K = g^{ab} \bmod p$, which Eve cannot
        compute (discrete logarithm).
\end{itemize}
\textbf{Open question}: how do Alice and Bob know they are talking to the right person,
and not to Eve in the middle (MITM)?
\end{rembox}

\section{Problem I: Authenticity}

\begin{defnbox}[title=Authenticity]
How do we \textbf{verify the identity} of the person to whom we are talking, while
avoiding Eve in the middle? Bob wants to talk to Alice; Alice wants to make sure she is
talking to Bob.
\end{defnbox}

\begin{defnbox}[title=Public-key cryptography: key pairs and PKI]
Each user has a \textbf{key pair}:
\begin{itemize}
  \item a \textbf{secret key} $SK$ (kept private), e.g.\ Bob's $SK = b$;
  \item a \textbf{public key} $PK$ (published), e.g.\ Bob's $PK = B = g^b$.
\end{itemize}
Public keys are published in a \textbf{Public Key Infrastructure (PKI)}, from which
anyone can fetch the public key of the person they want to talk to.
\end{defnbox}

\begin{rembox}[title=Link with Diffie-Hellman]
A DH public value $B = g^b$ \emph{is} a public key, and $b$ the secret key. The PKI
(e.g.\ Certificate Authorities, Lecture 5) is what binds a public key to an
\textbf{identity}: it answers ``is this really Alice's key?''.
\end{rembox}

\begin{defnbox}[title=Public-key encryption (confidentiality)]
To send a confidential message to Alice, Bob fetches Alice's \textbf{public key} from the
PKI and encrypts:
\[
  C = E_{PK_{\text{Alice}}}(M), \qquad M = D_{SK_{\text{Alice}}}(C).
\]
Only Alice, who holds $SK_{\text{Alice}}$, can decrypt. Eve also knows $PK$, but trying to
decrypt with the public key gives garbage. Example: \textbf{RSA-OAEP}.
\textbf{Don't design your own!}
\end{defnbox}

\begin{rembox}[title=Intuition: the padlock]
The public key is an \textbf{open padlock} that Alice hands out to everybody; the secret
key is the only key that opens it. Anyone can lock a box for Alice, only Alice can open
it.
\end{rembox}

\begin{defnbox}[title=Digital signatures (signing and verification)]
Bob \textbf{signs} with his \textbf{secret} key, Alice \textbf{verifies} with Bob's
\textbf{public} key (from the PKI):
\[
  S = \mathrm{Sign}_{SK_{\text{Bob}}}(M), \qquad
  V_{PK_{\text{Bob}}}(M, S) = \checkmark .
\]
\textbf{Unforgeability}: Eve cannot forge a signature $S$ on $M$ that verifies without
knowing $SK$. Example: \textbf{NIST DSA}. \textbf{Don't design your own!}
\end{defnbox}

\begin{rembox}[title=Encryption vs.\ signature: who uses which key?]
\begin{center}
\begin{tabular}{l|c|c}
                   & \textbf{uses $PK$} & \textbf{uses $SK$} \\ \hline
Encryption         & sender encrypts    & receiver decrypts  \\
Signature          & receiver verifies  & sender signs
\end{tabular}
\end{center}
Encryption: \emph{anyone} can write to Alice, only Alice can read. Signature: only Bob can
write (sign), \emph{anyone} can check it was Bob.
\end{rembox}

\begin{propbox}[title=Sign the hash of messages]
Public-key operations are \textbf{computationally costly} compared with most symmetric
key algorithms of equivalent security: signing and encrypting \textbf{is slow}. \\
$\Rightarrow$ \textbf{Sign the hash} of messages instead of the message itself:
\[
  H = h(M), \qquad S = \mathrm{Sign}_{SK}(H), \qquad V_{PK}(H, S) = \checkmark ,
\]
where Alice recomputes $H = h(M)$ from the received $M$.
\end{propbox}

\begin{defnbox}[title=Hash function]
A function $h(\cdot)$ that maps a message $M$ of \textbf{any length} to a \textbf{fixed
short-length} output $H = h(M)$. Three security properties:
\begin{itemize}
  \item \textbf{Pre-image resistance} --- given $H = h(M)$, it is difficult to find $M$.
  \item \textbf{Second pre-image resistance} --- given $M$, it is difficult to find an
        $M' \neq M$ such that $h(M') = h(M)$.
  \item \textbf{Collision resistance} --- it is difficult to find \emph{any} $M, M'$
        (with $M \neq M'$) such that $h(M) = h(M')$.
\end{itemize}
\end{defnbox}

\begin{rembox}[title=Understanding the three properties]
\begin{itemize}
  \item \textbf{Pre-image}: the hash is a one-way ``fingerprint'': from the fingerprint
        you cannot rebuild the person.
  \item \textbf{Second pre-image}: the attacker is given a \emph{fixed} $M$ and must
        match its hash.
  \item \textbf{Collision}: the attacker chooses \emph{both} messages freely --- an easier
        task, so collision resistance is the strongest requirement. By the birthday
        paradox, an $n$-bit hash only gives about $n/2$ bits of collision resistance.
\end{itemize}
Collisions must exist (infinitely many inputs, finitely many outputs); they must just be
\textbf{hard to find}.
\end{rembox}

\begin{propbox}[title=Standard hash functions]
\begin{itemize}
  \item \textbf{MD5} (1991): 128-bit hash --- \textbf{insecure}.
  \item \textbf{SHA-0, SHA-1}: 160 bits --- \textbf{insecure}.
  \item \textbf{SHA-2} (224/256/384/512) --- OK but slow.
  \item \textbf{SHA-3} (224/256/384/512).
\end{itemize}
\textbf{Don't design your own!} \\
\textbf{Uses}: support digital signatures, build HMAC, password storage, file integrity,
secure commitments, secure logging, blockchains, \ldots
\end{propbox}

\begin{propbox}[title=Which hash properties do signatures need?]
When signing $h(M)$, \textbf{pre-image resistance is not needed} ($M$ is sent anyway).
What matters:
\begin{itemize}
  \item \textbf{Second pre-image resistance} --- otherwise Eve finds $M' \neq M$ with
        $h(M') = h(M)$, and Bob's signature on $M$ is also valid on $M'$.
  \item \textbf{Collision resistance} --- otherwise Eve prepares two messages with the same
        hash, gets Bob to sign the harmless one, and reuses the signature on the other.
\end{itemize}
\end{propbox}

\begin{defnbox}[title={All together: asymmetric confidentiality and authentication}]
Users have \textbf{two pairs} of keys (2 secret keys, 2 public keys): an
\textbf{encryption} pair and a \textbf{signing} pair.
\begin{itemize}
  \item \textbf{Confidentiality}: $D_{SK}(C) = M$ with $C = E_{PK}(M)$.
  \item \textbf{Authentication}: $V_{PK}(M, S) = \checkmark$ with
        $S = \mathrm{Sign}_{SK}(M)$.
\end{itemize}
Bob sends to Alice
\[
  E_{PK_{\text{Alice}}}(M), \quad \mathrm{Sign}_{SK_{\text{Bob}}}\big(h(M)\big),
\]
encrypting with Alice's \textbf{encryption public key} and signing with his own
\textbf{signing secret key}; Alice decrypts with her secret key and verifies with Bob's
\textbf{verification public key}.
\end{defnbox}

\section{Problem II: Integrity}

\begin{defnbox}[title=Integrity]
Information cannot be \textbf{modified} by unauthorized parties. \\
Eve now \textbf{writes} on the channel: Bob sends ``Banana'', Eve tampers with the
ciphertext, and Alice decrypts ``Apple'' without noticing (recall the bit-flipping attack
on the OTP, Lecture 5).
\end{defnbox}

\begin{propbox}[title=Message integrity through digital signatures]
Digital signatures provide:
\begin{itemize}
  \item \textbf{Sender authenticity} --- the message comes from Bob;
  \item \textbf{Message integrity} --- the signature $S$ only verifies if the message $M$
        has not been modified;
  \item \textbf{Non-repudiation} --- Bob cannot later deny having signed $M$: only his
        secret key could produce $S$, and \emph{anyone} (a judge, a third party) can
        check it with his public key.
\end{itemize}
\end{propbox}

\subsection{Symmetric cryptography: block ciphers}

\begin{defnbox}[title=Block cipher]
And for symmetric cryptography? A block cipher encrypts a block with a pre-shared key:
$C_i = E_K(M_i)$, $M_i = D_K(C_i)$. The algorithms work on \textbf{blocks} the size of the
key, typically \textbf{128/256 bits}. \\
\textbf{Problem}: messages are longer than a block! $\Rightarrow$ requires iteration
$\Rightarrow$ block ciphers have a \textbf{mode of operation}.
\end{defnbox}

\begin{defnbox}[title={Mode 1: Electronic Code Book (ECB)}]
Straightforward scheme: split $M = M_1 M_2 M_3 M_4$ into blocks (each the size of the key)
and encrypt and decrypt \textbf{each block independently}:
\[
  C_i = E_K(M_i), \qquad M_i = D_K(C_i).
\]
\end{defnbox}

\begin{warnbox}[title=The problem with ECB]
If $M_i = M_j$ then $C_i = C_j$! Identical plaintext blocks give identical ciphertext
blocks: \textbf{large information leakage} due to the lack of randomness across blocks'
encryption (the same issue as substitution ciphers, but on whole blocks). Classic
illustration: an image encrypted with ECB still shows its outline.
\end{warnbox}

\begin{defnbox}[title={Mode 2: Cipher Block Chaining (CBC)}]
Add an \textbf{IV} and \textbf{propagate information across blocks} to introduce
randomness.
\begin{align*}
  \text{Encryption: } & C_0 = IV, \quad C_i = E_K(M_i \oplus C_{i-1}) \\
  \text{Decryption: } & C_0 = IV, \quad M_i = D_K(C_i) \oplus C_{i-1}
\end{align*}
\end{defnbox}

\begin{rembox}[title=Questions from the lecture]
\begin{itemize}
  \item What if the IV is incorrect? Is the full decryption wrong?
  \item Can you decrypt a block alone? What do you need?
\end{itemize}
(Hint: look at which values appear in the decryption formula of $M_i$.)
\end{rembox}

\begin{defnbox}[title={Mode 3: Counter mode (CTR)}]
Use an \textbf{increasing nonce} to add randomness \textbf{without dependencies between
blocks}. A \textbf{nonce} is a random number used only \textbf{once}.
\[
  \text{Encryption: } C_i = E_K(\mathit{Nonce} + i) \oplus M_i;
  \qquad
  \text{Decryption: } M_i = E_K(\mathit{Nonce} + i) \oplus C_i .
\]
\end{defnbox}

\begin{rembox}[title=CTR turns a block cipher into a stream cipher]
$E_K(\mathit{Nonce}+1), E_K(\mathit{Nonce}+2), \ldots$ is a pseudo-random \textbf{key
stream} XOR-ed with the message, exactly like a stream cipher (Lecture 5). Hence
decryption only needs $E_K$ (no $D_K$, answering the slide's question ``do we need a
decryption algorithm?''), blocks can be processed \textbf{in parallel}, and --- as with
any pad --- a nonce must \textbf{never be reused} with the same key.
\end{rembox}

\begin{propbox}[title=Block ciphers: strengths and weaknesses]
\textbf{Strengths}:
\begin{itemize}
  \item \textbf{High diffusion} --- information from one plaintext symbol is diffused
        into several ciphertext symbols;
  \item \textbf{Immunity to tampering} --- difficult to insert symbols without detection.
\end{itemize}
\textbf{Weaknesses}:
\begin{itemize}
  \item \textbf{Slow} --- an entire block must be accumulated before encryption /
        decryption can begin;
  \item \textbf{Error propagation} --- in some modes of operation errors affect several
        bits/blocks.
\end{itemize}
Individual modes may not fall in these categories and may offer a different trade-off.
It is the opposite profile of stream ciphers (Lecture 5).
\end{propbox}

\begin{keybox}[title=Modes of operation: summary]
\begin{itemize}
  \item \textbf{ECB} --- directly encrypt and decrypt single blocks. Large information
        leakage due to lack of randomness across blocks' encryption.
  \item \textbf{CBC} --- avoids ECB problems by using previous blocks to add randomness to
        every encryption. Propagates errors and prevents single-block decryption.
  \item \textbf{CTR} --- uses an increasing ``nonce'' to introduce randomness across
        blocks' encryptions. Enables parallel encryption and decryption.
\end{itemize}
\end{keybox}

\subsection{Integrity from symmetric cryptography: MACs}

\begin{defnbox}[title=Message Authentication Code (MAC)]
Takes a \textbf{fixed (short) size key} $K$ and a message $M$, and outputs a short
$MAC_K(M)$. It is \textbf{hard to generate} a pair $\{M, MAC_K(M)\}$ \textbf{without
knowing $K$}. Bob sends $\{M; MAC_K(M)\}$ using the pre-shared key $K$. Two properties:
\begin{enumerate}
  \item $M$ cannot be tampered with;
  \item the message comes from Bob/Alice (whoever holds $K$).
\end{enumerate}
$\Rightarrow$ \textbf{Mutual authentication BUT no sender authenticity towards a 3rd
party.}
\end{defnbox}

\begin{defnbox}[title=Repudiation]
With a MAC, both Alice and Bob hold the same key $K$, so \textbf{both} can produce
$MAC_K(M)$. Given $\{M = \text{``I don't like Charlie''}; MAC_K(M)\}$, Bob says ``it was
Alice'' and Alice says ``it was Bob'': a third party (Charlie) \textbf{cannot tell} who
wrote it.
\end{defnbox}

\begin{keybox}[title=MAC vs.\ digital signature]
\begin{center}
\begin{tabular}{l|c|c}
                         & \textbf{MAC}       & \textbf{Signature}  \\ \hline
Keys                     & one shared secret  & key pair $(SK, PK)$ \\
Integrity                & yes                & yes                 \\
Authenticity (to peer)   & yes                & yes                 \\
Convince a third party   & no                 & yes                 \\
Non-repudiation          & no                 & yes                 \\
Speed                    & fast               & slow
\end{tabular}
\end{center}
\end{keybox}

\begin{defnbox}[title=CBC-MAC]
Turning a block cipher into a MAC:
\[
  C_0 = 0 \ \text{[any fixed IV]}, \qquad C_i = E_K(M_i \oplus C_{i-1}), \qquad
  MAC_K(M_1 \ldots M_n) = C_n .
\]
Differences with respect to CBC encryption:
\begin{itemize}
  \item CBC-MAC is \textbf{deterministic} (fixed IV);
  \item only the \textbf{last block} is output.
\end{itemize}
\textbf{Limitation}: only secure if the length of message $M$ is known (fixed).
\end{defnbox}

\begin{rembox}[title=Why a fixed IV and why only the last block?]
A MAC must be \textbf{recomputable} by the verifier: with a random IV, Alice could not
recompute the same tag (and an attacker could adjust the IV to fix the first block).
Since every block is chained into the next, the last block depends on the whole message:
it is a short ``keyed fingerprint'' of $M$.
\end{rembox}

\subsection{All together: confidentiality and integrity (symmetric)}

\begin{defnbox}[title=Encrypt-and-MAC]
Send $E_K(M) \,\|\, MAC_K(M)$: encrypt the plaintext and MAC the \textbf{plaintext},
in parallel.
\begin{itemize}
  \item[$\times$] \textbf{No integrity on the ciphertext} $\to$ the cipher can be
        attacked; need to decrypt to know if valid.
  \item[$\checkmark$] Integrity of the plaintext can be verified.
  \item[$\times$] \textbf{May reveal information about the plaintext} $\to$ repeated
        messages give repeated MACs (the IV of the MAC is fixed; can be solved with a
        counter).
\end{itemize}
\end{defnbox}

\begin{defnbox}[title=MAC-then-Encrypt]
Send $E_K\big(M \,\|\, MAC_K(M)\big)$: MAC the plaintext, then encrypt plaintext and MAC
together.
\begin{itemize}
  \item[$\times$] \textbf{No integrity of the ciphertext}: (in theory) possible to change
        the ciphertext and still have a valid MAC; need to decrypt to know if valid.
  \item[$\checkmark$] Integrity of the plaintext can be verified.
  \item[$\checkmark$] No information on the plaintext either, since it is encrypted.
\end{itemize}
\end{defnbox}

\begin{defnbox}[title=Encrypt-then-MAC]
Send $C \,\|\, MAC_K(C)$ with $C = E_K(M)$: encrypt first, then MAC the
\textbf{ciphertext}.
\begin{itemize}
  \item[$\checkmark$] \textbf{Integrity of the ciphertext} $\to$ ensures you only read
        valid messages! The cipher cannot be attacked.
  \item[$\checkmark$] Integrity of the plaintext can be verified.
  \item[$\checkmark$] No information on the plaintext either, since it is encrypted.
\end{itemize}
\end{defnbox}

\begin{keybox}[title=Which composition to use?]
\begin{center}
\begin{tabular}{l|c|c|c}
                     & \textbf{E\&M} & \textbf{MtE} & \textbf{EtM} \\ \hline
Ciphertext integrity & $\times$      & $\times$     & $\checkmark$ \\
Plaintext integrity  & $\checkmark$  & $\checkmark$ & $\checkmark$ \\
Hides the plaintext  & $\times$      & $\checkmark$ & $\checkmark$
\end{tabular}
\end{center}
\textbf{Encrypt-then-MAC} wins: the receiver checks the MAC \emph{before} decrypting, so
forged ciphertexts are rejected without ever being fed to the decryption algorithm
(fail-safe default, Lecture 2). Use separate keys for encryption and MAC.
\end{keybox}

\begin{keybox}[title={Summary: what we have learned about applied cryptography}]
\begin{itemize}
  \item \textbf{Symmetric cryptography}
        \begin{itemize}
          \item Confidentiality: stream ciphers, block ciphers (\textbf{modes of
                operation}!).
          \item Integrity / authentication: \textbf{Message Authentication Codes (MACs)}.
        \end{itemize}
  \item \textbf{Asymmetric cryptography}
        \begin{itemize}
          \item Confidentiality: \textbf{encryption} with the receiver's public key.
          \item Integrity / authentication: \textbf{digital signatures} (+
                non-repudiation).
        \end{itemize}
  \item \textbf{Hash functions}
        \begin{itemize}
          \item Three security properties (pre-image, second pre-image, collision).
          \item Support digital signatures + other functions.
        \end{itemize}
\end{itemize}
\end{keybox}
